\begin{abstract}
This article studies how 2-gaps evolve under successive prime filters and
sharpens the boundary between complete-period survival and square-window
placement. Complete-period CRT arguments prove that 2-gaps persist
globally: an incoming odd prime removes exactly two copy classes of each
old 2-gap. Those counts do not force a survivor into a particular
interval $[Q,Q^2)$.

The full argument reduces survival to a single signed quantity. A
weighted conservation law and its Cauchy--Schwarz corollary give a sharp
terminal threshold on the harmful-excess energy. A closed exhaustion
argument---capacity envelopes, native-period Bessel, fixed and moving
cuts, and the stability-gap repair---shows that no unsigned-capacity
route can clear that threshold. The new contribution is consequently
signed and local. At filter $7$, exact residue order gives the sharp
interval bound $|b_7|\le18/7$, replacing a capacity estimate that grows
quadratically with the window scale. More generally, when an old period
is copied through an incoming prime $r$, the centered harmful excess in
copy block $j$ is exactly $B_j=d_t+d_{t-2}$ for two entries of the
centered old-start histogram modulo $r$. Consequently,
\begin{equation*}
\sum_{j=0}^{r-1}B_j^2
=2V_r+2\sum_{t\bmod r}d_td_{t-2}
\le4V_r.
\end{equation*}

This composes residue energy with the weighted harmful-excess survival
criterion over complete old-period blocks. It does not control the two
partial boundary fragments of an arbitrary interval, nor does it prove a
relative residue-energy estimate across an unbounded filter chain. Those
are now the precise remaining arithmetic obligations, framed by a scale
conflict that limits fixed-seed averaging and a Type-II barrier that any
prime-producing sieve must overcome.
\end{abstract}
